\documentclass[12pt]{article}
\usepackage[utf8]{inputenc}
\usepackage[T1]{fontenc}
\usepackage{hyperref}
\title{Bena Language Resources for K-12 Curriculum Development}
\author{Compiled by Claude Sonnet 5.5 (claude-sonnet-5-5) for Andr\'as Kornai}
\date{}
\begin{document}
\maketitle

\section{Introduction}
Bena of Nigeria (ISO 639-3 yun; Glottocode bena1260), also called Yungur or Yongor, is an Adamawa language of Adamawa State, Nigeria, with about 95,000 speakers (1992 figure cited on Wikipedia). It must not be confused with Bena (bez) of Tanzania, Benabena (bef) of Papua New Guinea, or the village of Bena in Kebbi State. The searches found almost no educational material for this language.

\section{Primary Educational Resources}
None found.

\section{Linguistic and Reference Materials}
\begin{itemize}
\item Wikipedia article on Bena (Adamawa), which cites Ethnologue and Blench's \emph{Atlas of Nigerian Languages}: \url{https://en.wikipedia.org/wiki/Bena_language_(Adamawa)}
\item Wikipedia disambiguation page: \url{https://en.wikipedia.org/wiki/Bena_language}
\item ASJP word-list entries named Bena (the entries were not checked to confirm they are the Nigerian language): \url{https://asjp.clld.org/languages/BENA} and \url{https://asjp.clld.org/languages/BENA_2}
\end{itemize}

\section{Religious and Community Resources}
None found.

\section{Digital and Audio Resources}
None found.

\section{Curriculum Development Context}
Practically no teaching or documentation material was found. A curriculum would first require basic documentation, orthography work and community partnership. Roger Blench's atlas and Ethnologue (subscription) are cited by Wikipedia as further leads but were not seen directly.

\end{document}